% Section 10: Non-faithfulness: the residual catalogue
\label{sec:residual-body}

The previous sections produced no-go results one at a time.  This section collects them,
because together they say something stronger than each does separately: there is a
\emph{systematic} pattern in which quantity fails to determine construction, and naming
the pattern is more useful than proving its instances.  The section also records the
positive counterweights, since a catalogue of failures without a statement of what a
faithful carrier looks like would be misleading.

\subsection{The pattern}

\begin{proposal}[The residual pattern]\label{prop:residual-pattern}
Let $\sigma$ be a quantity computed from a process by an observer $Q$.  In every
instance catalogued below, $\sigma$ is a \emph{partial quotient}: it is invariant under a
group of construction changes that the process itself does not admit.  The correct
response is never to replace $\sigma$ by a finer scalar; it is to retain a residual
recording the orbit of the forgotten changes.
\end{proposal}

\begin{remark}
Proposition~\ref{prop:residual-pattern} is a \textsf{Proposal} because "partial quotient"
is a description, not a hypothesis with a theorem attached.  Its content is
methodological and its justification is the list: in each case below, the residual is
exactly what repairs the deficiency, and no case is repaired by refining the scalar.
\end{remark}

\subsection{Spectral shadows}

\begin{theorem}[Spectrum does not determine history]\label{thm:spectrum-history}
Let $x$ be a single input.  The programs $\mathrm{add}(\mathrm{copy}(x))$ and
$\mathrm{scale}_{2}(x)$ agree on realized value, on derivative, on all iterates, and on
the action spectrum, which is $\{2\}$ in both cases; they differ in the construction
graph, in the occurrence structure, and in the source history.  Hence no observer that
sees only values, derivatives, iterates, and spectra can separate them.
\end{theorem}

\begin{proof}
Both programs induce the same endomorphism $x\mapsto 2x$ on the declared state, so every
functional of the induced action agrees.  They differ in the number of occurrences of
the source and in the recorded lineage, which are not functions of the induced action.
\end{proof}

\begin{proposition}[The affine carrier is not diagonalisable into a spectrum]
\label{prop:affine-spectrum}
In the affine carrier, scaling acts diagonally with spectrum
$\{m^{-1},\dots,m^{-K}\}$, while addition acts unipotently and is nilpotent on the
relevant graded pieces.  Consequently a bare eigenvalue set cannot encode addition: the
characteristic polynomial of a unipotent operator is $(\lambda-1)^{d}$ regardless of the
translation amount.
\end{proposition}

\begin{proof}
Scaling conjugates $A\mapsto M AM^{-1}$ and stretches the translation grid by $m$ at each
graded level, giving the stated spectrum.  Addition is the exponential of a nilpotent
translation, whose characteristic polynomial carries no information about the amount.
\end{proof}

\subsection{Two faithful carriers, for contrast}

\begin{theorem}[Marked second-jet faithfulness]\label{thm:second-jet}
The positive affine fragment embeds faithfully into the marked second jet at the
projective point at infinity,
\[
  m=c_1^{-1},\qquad a=-c_2c_1^{-2},
\]
where $(c_1,c_2)$ are the marked jet coefficients.  The faithfulness is a statement about
\emph{marked} jets: the unmarked jet does not suffice, since it is defined only up to the
chart that the marking fixes.
\end{theorem}

\begin{proof}
Composition of affine maps corresponds to multiplication in
$\R[[q]]/(q^{3})$ on the marked coefficient pair; solving for $(m,a)$ gives the displayed
formulas, and the correspondence is a bijection onto the positive affine fragment within
the marked class.
\end{proof}

\begin{theorem}[Co-generation of polynomial and matrix readings]\label{thm:cogeneration}
Transport is homogeneous of degree one in declared symbolic probe generators, and the
extracted coefficient table equals the transpose of the forward Jacobian.  Hence the
polynomial-like and matrix-like presentations of one process are co-generated, rather
than independently imposed.
\end{theorem}

\begin{proof}[Status]
Proved for the declared probe class by linearity of the demand section.  The phrase
"after a chart" is essential: the identification of the coefficient table with the
transpose requires the port chart to be declared, and without it the two readings are
merely isomorphic.
\end{proof}

\begin{remark}[The moral of the contrast]
Theorems~\ref{thm:second-jet} and~\ref{thm:cogeneration} are the positive half of the
catalogue.  Both work by making the marking or the chart explicit.  The failures below
are exactly the cases where a marking or chart was silently omitted, and the residual
records what the omission cost.
\end{remark}

\subsection{The observation tower is strict}

\begin{theorem}[Strict observation tower]\label{thm:observation-tower}
The three observation levels
\[
  \cO_{\mathrm{prog}}\ \supsetneq\ \cO_{\mathrm{or}}\ \supsetneq\ \cO_{\mathrm{proj}}
\]
are strictly separated: projective spectral data is strictly coarser than oriented action
data, which is strictly coarser than program-aware data.  Each strictness is witnessed by
a pair of programs identified at the lower level and separated at the higher.
\end{theorem}

\begin{proof}
The projective level forgets the central sign
(Proposition~\ref{prop:order-two-four}), and the oriented level forgets occurrence
structure (Theorem~\ref{thm:spectrum-history}).  Both forgettings are non-injective on the
finitary fixtures exhibited there.
\end{proof}

\subsection{Objectification and effective completion}

\begin{nogo}[Objectification is not natural]\label{nogo:objectification}
Let $\mathrm{Obj}$ send a finite cut presentation to its simplest surreal value.  Then
$\mathrm{Obj}$ is non-injective, and it is not natural:
\[
  g\bigl(\mathrm{Obj}\,C\bigr)\neq \mathrm{Obj}\bigl(gC\bigr)
  \qquad\text{for } g(z)=3z \text{ and a suitable } C .
\]
Consequently no analysis factoring through the objectified value admits a left-inverse
synthesis.
\end{nogo}

\begin{proof}
Non-injectivity is immediate, since distinct presentations may share a simplest value.
For non-naturality, the two sides of the displayed inequality are computed on an explicit
finite $C$; the objectified value is simplified before scaling on the left and after
scaling on the right, and simplicity is not preserved by scaling.
\end{proof}

\begin{proposition}[Effective one-sidedness]\label{prop:effective-onesided}
For a universal prefix-free machine $U$, let
\[
  \Omega_U=\sum_{U(p)\downarrow}2^{-|p|}\le 1
\]
be its halting mass, and let $(L_\Omega\mid R_\Omega)$ be the associated cut.  Then
$L_\Omega$ is computably enumerable while $R_\Omega$ is not.  Hence a computational cut
is intrinsically asymmetric.
\end{proposition}

\begin{proof}[Status]
The Kraft bound is a standard measure-theoretic computation on the cylinder sets of a
prefix-free set.  The asymmetry is the standard halting asymmetry: membership on the left
is witnessed by a halting computation, while membership on the right would decide
non-halting.
\end{proof}

\begin{proposition}[Mass cannot repair the fibre]\label{prop:mass-not-fibre}
The objectification fibre mass
\[
  m_{U,Q}(x)=\sum_{\mathrm{Obj}_Q(U(p))=x}2^{-|p|}
\]
is a lower-semicomputable semimeasure.  It quantifies the fibre but cannot recover its
members.
\end{proposition}

\begin{proof}
Lower semicomputability follows from the definition as a sum over enumerated halting
computations.  Recovery fails because a sum over a set determines neither its cardinality
nor its elements.
\end{proof}

\begin{proposition}[No total finite coding of $\No$]\label{prop:no-finite-coding}
There is no set-sized family of finite codes surjecting onto the surreal numbers.
\end{proposition}

\begin{proof}
A set-sized family of finite codes is a set, and $\No$ is a proper class; no set surjects
onto a proper class in the relevant setting.
\end{proof}

\subsection{Snapshots, scopes, and traces}

\begin{proposition}[Equal snapshots, nonempty interval]\label{prop:snapshot-not-empty}
For exact program slices between nested cuts, equal frontier snapshots do not imply an
empty interval: an explicit discard witnesses two distinct slices with the same snapshot.
\end{proposition}

\begin{proposition}[Zero-event incidence is not single-valued]\label{prop:zero-event-fibre}
A zero-event graft frame is compatible with two distinct identity slices fixed by an
independent surgery.  Hence the available frame and slice data admit no canonical
single-valued placement of scope in cut; the fibre has two elements.
\end{proposition}

\begin{proposition}[Trace arithmetic does not identify the sides]\label{prop:trace-sides}
For the three-sided trace calibration, the temporal and spatial residuals vanish while
the constructional residue $(+1,-1,0,0,0)$ does not.  Consequently the two sides are not
semantically identified by their trace arithmetic.
\end{proposition}

\begin{proposition}[Count factorization and the recovery obstruction]
\label{prop:count-factorization}
Let $T$ be the trace count and $n$ the construction length.  Then
\[
  T(p)=\bigl(n,\ n-1,\ 3(n-1)\bigr),
\]
a positive factorization through construction length.  However the projection
$(T,S)\mapsto C$ is not single-valued: an explicit pair of constructions shares $(T,S)$
and differs in $C$.
\end{proposition}

\begin{proof}[Status]
The factorization is an exact finite count.  The obstruction is an explicit pair, checked
on the recorded data; it is a counterexample to single-valued recovery, not a general
classification of the fibres of $(T,S)\mapsto C$.
\end{proof}

\begin{proposition}[Period covectors do not determine history]\label{prop:period-not-history}
A period covector is a functional of a homology class and therefore constant on the
construction paths that realise that class.  It cannot reconstruct the construction
history.
\end{proposition}

\subsection{Minimal expressive repair}

One result in the catalogue deserves separate billing, because it gives an exact
\emph{minimum} rather than only an obstruction.

\begin{theorem}[Order-embedding obstruction]\label{thm:order-obstruction}
Let a coarse observation erase an incomparability, i.e.\ two elements become
indistinguishable in the observation order.  Then no single added scalar coordinate can
restore the lost order, whereas two monotone Boolean coordinates can, and two is the
minimum.
\end{theorem}

\begin{proof}
A single scalar added to a coarse observation produces a total preorder on the fibre,
since scalars are totally ordered; a total preorder cannot contain an incomparability.  Two
monotone Boolean coordinates separate the four possible cells, which suffices to embed the
erased relation, and one coordinate gives only two cells.  Hence two is both sufficient
and necessary.
\end{proof}

\begin{remark}[Why this is the catalogue's high point]
Theorem~\ref{thm:order-obstruction} converts "coarse observation loses a relation" from a
complaint into a quantitative statement with an exact constant.  It is the model for what
the rest of the catalogue would need in order to become equally informative.
\end{remark}

\subsection{Scale and winding}

\begin{proposition}[Scale forgetting]\label{prop:scale-forgetting}
Let $\lambda=\varepsilon Lz$ with $L=(2I+C)/3$ and $C$ the cyclic permutation matrix,
$C^{3}=I$.  Then $L$ is invertible over $\Q$, with
\[
  (2I+C)\bigl(4I-2C+C^{2}\bigr)=9I ,
\]
so the scale-marked triple is recoverable provided $\varepsilon\neq0$.  If the scale
marker is discarded, the pairs $(\varepsilon,z)$ and $(\varepsilon/2,2z)$ share $\lambda$,
and recovery fails.
\end{proposition}

\begin{proof}
Expanding $(2I+C)(4I-2C+C^{2})=8I-4C+2C^{2}+4C-2C^{2}+C^{3}=9I$ using $C^{3}=I$ gives the
inverse, hence invertibility.  The collision is immediate from
$\lambda=(\varepsilon/2)L(2z)$.
\end{proof}

\begin{nogo}[Winding counts do not determine the global action]\label{nogo:winding}
Abelian winding counts do not determine the global action on
$\PP^{1}\setminus\{0,1,\infty\}$: there are words of total winding zero whose commutator is
nontrivial,
\[
  ABA^{-1}B^{-1}=\begin{pmatrix}13&8\\ 8&5\end{pmatrix}\neq I .
\]
Moreover endpoint closure does not imply identity of the transported homology: a diamond
path of winding $+1$ transports $(0,1)$ to $(2,1)$.
\end{nogo}

\begin{proof}
Two computations on explicit words.  The first exhibits a nontrivial commutator with zero
abelian shadow, so the abelianisation is not faithful.  The second exhibits a closed
endpoint with nontrivial transport.
\end{proof}

\subsection{The catalogue}

\begin{center}
\small
\begin{longtable}{@{}p{0.30\textwidth}p{0.56\textwidth}@{}}
\toprule
Shadow & What it fails to determine \\
\midrule
\endhead
value, derivative, iterates, spectrum & construction graph, occurrence, source history
  (Thm~\ref{thm:spectrum-history}) \\
eigenvalue set & the translation amount of a unipotent factor
  (Prop~\ref{prop:affine-spectrum}) \\
projective spectral data & the central sign; oriented action
  (Thm~\ref{thm:observation-tower}) \\
oriented action data & occurrence structure
  (Thm~\ref{thm:observation-tower}) \\
objectified surreal value & the presentation, and it is not natural
  (No-go~\ref{nogo:objectification}) \\
fibre mass of an objectification & the members of the fibre
  (Prop~\ref{prop:mass-not-fibre}) \\
frontier snapshot & emptiness of an interval
  (Prop~\ref{prop:snapshot-not-empty}) \\
frame and slice data & a single-valued scope placement
  (Prop~\ref{prop:zero-event-fibre}) \\
trace arithmetic & identification of the two sides
  (Prop~\ref{prop:trace-sides}) \\
count pair $(T,S)$ & the construction $C$
  (Prop~\ref{prop:count-factorization}) \\
period covector & construction history
  (Prop~\ref{prop:period-not-history}) \\
discarded scale marker & the marked triple
  (Prop~\ref{prop:scale-forgetting}) \\
abelian winding & the global action
  (No-go~\ref{nogo:winding}) \\
boundary-return word & a periodic orbit
  (No-go~\ref{nogo:return-not-period}) \\
directional energy & any nontrivial characteristic
  (No-go~\ref{nogo:zero-energy-trap}) \\
\bottomrule
\end{longtable}
\end{center}

\subsection{What the catalogue implies}

\begin{enumerate}
\item The failures are not ad hoc.  In each row the shadow is invariant under changes the
  process does not admit, and the repair is a residual rather than a finer scalar
  (Proposal~\ref{prop:residual-pattern}).
\item Where a residual is supplied, the theory works: marked jets and declared port charts
  give genuine faithfulness (Theorems~\ref{thm:second-jet}, \ref{thm:cogeneration}).
\item At least one obstruction has an exact minimum
  (Theorem~\ref{thm:order-obstruction}), which is the standard the rest should meet.
\item Two of the rows are about \emph{absence of structure} rather than loss of
  information: No-go~\ref{nogo:return-not-period} (no feedback operator) and
  No-go~\ref{nogo:zero-energy-trap} (no nontriviality constraint).  These are the
  invitations to Part~II, where the missing structures are supplied syntactically.
\end{enumerate}
